% Edited conversion of source pp. 26--27. Equations reconstructed from the rendered PDF.
\ReportHeading
  {Про стійкість коливань пружної пластини, яка поділяє ідеальні рідини різної густини в циліндричному резервуарі}
  {Ю.~М. Кононов, А.~А. Зарецький}
  {Інститут прикладної математики і механіки НАН України}
  {}

У лінійній постановці розглянуто гідропружну задачу про вільні коливання тонкої пластини, яка горизонтально розділяє ідеальні нестисливі рідини різної густини в жорсткому циліндричному резервуарі довільного поперечного перерізу. Для розв’язання неоднорідного бігармонічного рівняння використано фундаментальну систему розв’язків відповідного однорідного бігармонічного рівняння та власні форми коливань ідеальної рідини у циліндричній порожнині.

Отримано частотне рівняння для довільних способів закріплення контуру пластини. Для затиснутої пластини його спрощено шляхом розкладання однорідного бігармонічного рівняння на два гармонічні рівняння та застосування формули Гріна для оператора Лапласа. У цьому випадку частотне рівняння не залежить від фундаментальної системи розв’язків, яка сама залежить від невідомої частоти.

Власні частоти та форми сумісних коливань пластини і рідини визначаються із системи інтегро-диференціальних рівнянь і граничних умов
\begin{equation}
 \Delta^2 w-2P\Delta w+q
 =\frac{\omega^2}{D}\sum_{n=1}^{\infty}\frac{a_n w_n}{k_n}\psi_n+C,
\end{equation}
\begin{equation}
 w_n=\frac{1}{N_n^2}\int_S w\psi_n\,\dd S,
 \qquad
 \int_S w\,\dd S=0,
\end{equation}
\begin{equation}
 \left(L_p[w]\right)_{\Gamma}=0,
 \qquad p=1,2.
\end{equation}
Тут $w$ --- прогин пластини,
\[
 P=\frac{T}{2D},\qquad
 q=\frac{g\Delta\rho-m_0\omega^2}{D},\qquad
 a_n=\rho_1\coth\kappa_{1n}+\rho_2\coth\kappa_{2n},
\]
\[
 \Delta\rho=\rho_2-\rho_1,\qquad
 \kappa_{in}=k_n h_i,\qquad
 N_n^2=\int_S\psi_n^2\,\dd S.
\]
Функції $\psi_n$ і числа $k_n$ є власними функціями та відповідними власними числами коливань ідеальної рідини у циліндричній порожнині довільного поперечного перерізу; $C$ --- довільна стала; $\Gamma$ --- контур області $S$; $L_1$ і $L_2$ --- диференціальні оператори граничних умов закріплення пластини.

Для затиснутого контуру частотне рівняння має вигляд
\begin{equation}
 \sum_{n=1}^{\infty}
 \frac{k_n}{\omega^2 a_n-k_n d_n}
 \left(\frac{B_n}{N_n}\right)^2=0,
\end{equation}
де
\[
 B_n=\left.\psi_n\right|_{\Gamma}=\mathrm{const},
 \qquad
 d_n=(Dk_n^2+T)k_n^2+g\Delta\rho-m_0\omega^2.
\]
Рівняння (4) охоплює низку окремих випадків. Для мембрани слід покласти $D=0$, а за відсутності верхньої рідини --- $\rho_1=0$.

Для прямокутного каналу шириною $2a$ у випадку плоских коливань
\[
 \psi_n(x,y)=\psi_n(x)=\cos k_n(x+a),\qquad
 k_n=\frac{\pi n}{2a},\qquad
 B_n=(-1)^n,
 \qquad N_n^2=a.
\]
У цьому випадку рівняння~(4) збігається з рівнянням, отриманим у~\cite{fluidplate:kononov2017}.

Для кругової пластини та рідини у круговому циліндрі частотне рівняння набуває вигляду
\begin{equation}
 \sum_{n=1}^{\infty}
 \frac{\xi_{nm}^3}
 {(\xi_{nm}^2-m^2)\,[a\omega^2 a_{nm}-\xi_{nm}d_{nm}]}=0.
\end{equation}
Безрозмірний параметр $\xi_{nm}=k_{nm}a$ визначається з рівняння $J_m'(\xi_{nm})=0$, де $J_m$ --- функція Бесселя першого роду порядку $m$.

Для великих частот, $n\gg1$, з рівняння~(5) випливає наближена формула
\[
 \omega_n^2\approx
 \frac{\xi_{nm}^2\bigl[(Dk_{nm}^2+T)k_{nm}^2+g\Delta\rho\bigr]}
 {(a_{nm}+m_0k_{nm})a},
\]
а отже, наближена умова стійкості сумісних коливань кругової пластини і рідини
\begin{equation}
 (Dk_{1m}^2+T)k_{1m}^2+g\Delta\rho>0.
\end{equation}
Вона завжди виконується за $T\ge0$ і $\rho_2\ge\rho_1$.

Чисельні дослідження для невагомості та безмасової пластини, $g=0$, $m_0=0$, показали, що найбільші частоти відповідають вільному краю, дещо менші --- затиснутому, а найменші --- опертому. Коли глибини заповнення перевищують радіус циліндра, їх подальше збільшення майже не впливає на частоти. Зі зростанням густини верхньої рідини частоти зменшуються, що може призвести до втрати стійкості сумісних коливань.

Роботу підтримано грантом Simons Foundation SFI-PD-Ukraine-00017674 (Yuriy Kononov, Andrii Zaretskiy).

\begin{thebibliography}{9}
\bibitem{fluidplate:kononov2023}
Yu.~M. Kononov,
``On the solution of a complicated biharmonic equation in a hydroelasticity problem,''
\emph{Journal of Mathematical Sciences}, vol.~274, no.~3, pp.~340--352, 2023.
\url{https://doi.org/10.1007/s10958-023-06604-w}.

\bibitem{fluidplate:kononov2017}
Ю.~Н. Кононов, А.~А. Лимарь,
``О колебании прямоугольной пластины, разделяющей идеальные жидкости разной плотности в прямоугольном канале с одним упругим основанием,''
\emph{Проблеми обчислювальної механіки і міцності конструкцій}, вип.~26, с.~79--96, 2017.

\bibitem{fluidplate:kononov2025}
Ю.~М. Кононов, А.~А. Зарецький, С.~В. Сапунов,
``Про коливання кругової пластини, яка горизонтально розділяє рідини різної густини у циліндричному резервуарі,''
\emph{Праці ІПММ НАН України}, т.~40, №~2, с.~138--153, 2025.
\url{https://doi.org/10.37069/1683-4720-2025-39-12}.
\end{thebibliography}
